\section{The environment, which is also the verifier}
\label{sec:env}

\subsection{One object, two names}
\label{sec:env-one-object}

In this benchmark the environment and the verifier are the same object. Stepping it is the same
act as asking for a verdict: the model submits a candidate fold, and the reply is either the
resulting state or a refusal. There is no separate validity oracle, no second artifact that could
disagree with the first, and no learned component in the loop.

\subsection{The division of labour}
\label{sec:env-division}

\begin{table}[h]
\caption{Who does what. The environment does not search.}
\label{tab:division}
\begin{center}
\begin{tabular}{lp{0.62\textwidth}}
\multicolumn{1}{c}{\bf COMPONENT}  &\multicolumn{1}{c}{\bf RESPONSIBILITY} \\
\hline \\
Environment / verifier & State update and legality. It does not search and does not choose \\
Model                  & Proposes the next fold, decides where to go, recognises when to backtrack \\
\end{tabular}
\end{center}
\end{table}

This division is what makes the benchmark measure something. The obvious objection to any
tool-augmented setup is that the tool is doing the work, and here the answer is structural rather
than rhetorical: deciding global flat-foldability is NP hard, so the environment cannot search its
way to an answer even in principle. It can only say whether one proposed step is something paper
could do. Compressing an exponential search into a feasible number of queries is exactly what the
environment cannot supply and exactly what the model is being measured on. The referee is not a
player.

\subsection{Refusals are named, isn't that you shoulr mention a check, each step check and final step check, final step check now is missing , should mention here--------}
\label{sec:env-refusals}

An illegal fold returns a structured refusal with a name, not a boolean.

\begin{table}[h]
\caption{The refusal taxonomy. A rejection arrives already labelled with the class it violated.}
\label{tab:refusals}
\begin{center}
\begin{tabular}{lp{0.60\textwidth}}
\multicolumn{1}{c}{\bf REFUSAL}  &\multicolumn{1}{c}{\bf MEANING} \\
\hline \\
\texttt{would-tear} & The moving layers are joined to stationary paper away from the fold line, so
  the sheet would come apart. The binding constraint of the some-layers tier \\
\texttt{no-crease} & The fold moves layers without creasing anything, tearing them free of the
  sheet rather than folding them \\
\texttt{nothing-to-move} & The fold line misses the selected layers entirely \\
\texttt{direction-impossible} & A run taken from the bottom of the stack was asked to fold over,
  or the reverse \\
\texttt{bad-line} & The fold line is malformed, or names no side of itself \\
\end{tabular}
\end{center}
\end{table}

The names are the design. A model told \texttt{would-tear} can act on the information; a model told
``illegal'' cannot. This also turns error analysis into a taxonomy that exists before the
experiment: every rejection arrives already labelled with the class it violated, so a reject is a
diagnosis rather than a dead end.

FInal step check. - The four constraint
classes that Flat-Folder checks (taco-taco, taco-tortilla, tortilla-tortilla, transitivity)
\citep{akitaya-flatfolder,flatfolder} remain the right vocabulary for global terminal states; they
are not what a single legal step needs checking against here.

\subsection{What the model observes - you should add a screen shot of simulator software !!!!!!!!}
\label{sec:env-observes}

On success the environment returns the new state together with two views: a top-down X-ray of the
stack, and an exploded view of the layers. Every intermediate state in this tier is flat, so a
second camera angle would show the same silhouette rotated and would carry no new information;
what carries information is layer structure, which is why the second view is an explosion rather
than a rotation. Positions are reported as coordinates on the original sheet together with an
integer layer index. There is no continuous vertical coordinate at this tier, and a rendering that
looks three-dimensional is not evidence of one; thickness belongs to a different problem.

